\documentclass[11pt,a4paper]{article}
\usepackage[utf8]{inputenc}
\usepackage[T1]{fontenc}
\usepackage[margin=2.5cm]{geometry}
\usepackage{graphicx}
\usepackage{booktabs}
\usepackage{float}
\usepackage{hyperref}

\title{MAC5921 - Deep Learning — EP1}
\author{André Nogueira Ribeiro}
\date{\today}
\begin{document}
\maketitle

\section{Declaração sobre uso de IA}

O código implementado para este trabalho foi gerado com ajuda da IA, utilizando uma IDE com modelos integrados. Abaixo se encontra com mais detalhes como foi utilizada a ajuda desta ferramenta:

\begin{itemize}
\item Códigos associados às implementções de funções e orquestrações para rodar os experimentos foram feitos com muito cuidado, o que inclui revisões detalhadas das sugestões geradas por IA, discussões com os modelos sobre diferentes possibilidades e razões para certas sugestões.
\item Códigos de teste foram revisados de maneira um pouco mais superficial, conferindo se o que estava sendo testado fazia sentido e em alguns casos indo mais a fundo no como isso era feito.
\item Códigos para geração de figuras e tabelas com os resultados foram feitos com pouca intervenção da minha parte.
\item Este relatório foi feito sem nenhuma ajuda de modelos, exceto para pequenas formatações e sintáxe de \LaTeX.
\end{itemize}

\section{Código}

O código, incluindo este relatório feito em \LaTeX, pode ser encontrado em:
\begin{itemize}
\item \href{https://github.com/deconr/MAC5921-DeepLearning}{https://github.com/deconr/MAC5921-DeepLearning}.
\end{itemize}

\section{Planejamento dos experimentos}

\section{Descrição dos experimentos realizados}



\end{document}